% This is text associated with the open textbook "Open Electromagnetics"
% See immediately below for author, date
% This work is licensed under the Creative Commons Attribution-ShareAlike 4.0 International License. (CC BY-SA 4.0) 
% To view a copy of the license, visit http://creativecommons.org/licenses/by-sa/4.0/

%ORB TITLE Directivity and Gain
%ORB AUTHOR 0        # S.W. Ellingson, Virginia Tech (ellingson@vt.edu)
%ORB DATE 20191212   # yyyymmdd
%ORB ABSTRACT 

\label{m0203_Directivity_and_Gain} % <-- Must match name of this file and module directory name.  Don't mess with this!
\index{directivity|(}

A transmitting antenna does not radiate power uniformly in all directions.
Inevitably more power is radiated in some directions than others.
\emph{Directivity} quantifies this behavior.
In this section, we introduce the concept of directivity and the related concepts of \emph{maximum directivity} and \emph{antenna gain}.

Consider an antenna located at the origin.
The power radiated in a single direction $(\theta,\phi)$ is formally zero.
This is because a single direction corresponds to a solid angle of zero, which intercepts an area of zero at any given distance from the antenna.
Since the power flowing through any surface having zero area is zero, the power flowing in a single direction is formally zero.
Clearly we need a different metric of power in order to develop a sensible description of the spatial distribution of power flow.

The appropriate metric is \emph{spatial power density}; that is, power per unit area, having SI base units of W/m$^2$.
Therefore, directivity is defined in terms of spatial power density in a particular direction, as opposed to power in a particular direction.
Specifically, directivity in the direction $(\theta,\phi)$ is:
%
\begin{equation}
D(\theta,\phi) \triangleq \frac{S({\bf r})}{S_{ave}(r)}
\label{m0203_eDirDef}
\end{equation}
%
In this expression, $S({\bf r})$ is the power density at $(r,\theta,\phi)$; i.e., at a distance $r$ in the direction $(\theta,\phi)$. 
$S_{ave}(r)$ is the \emph{average} power density at that distance; that is, $S({\bf r})$ averaged over all possible directions at distance $r$.
Since directivity is a ratio of power densities, it is unitless.
Summarizing:
%%%%%%%%%%%%%%%%%%%%%%%%%%%%%%%%%%%%%%%%%%%%%%%
%ORB HIGHLIGHT
\begin{mdframed}[backgroundcolor=yellow!20,nobreak=true] 
Directivity is ratio of power density in a specified direction to the power density averaged over all directions at the same distance from the antenna.
\end{mdframed}
%%%%%%%%%%%%%%%%%%%%%%%%%%%%%%%%%%%%%%%%%%%%%%%

Despite Equation~\ref{m0203_eDirDef}, directivity does not depend on the distance from the antenna.  
To be specific, directivity is the same at every distance $r$. 
Even though the numerator and denominator of Equation~\ref{m0203_eDirDef} both vary with $r$, one finds that the distance dependence always cancels because power density and average power density are both proportional to $r^{-2}$.  
This is a key point: Directivity is a convenient way to characterize an antenna because it does not change with distance from the antenna.

In general, directivity is a function of direction.  However,
one is often not concerned about all directions, but rather only the directivity in the direction in which it is maximum.
In fact it is quite common to use the term ``directivity'' informally to refer to the maximum directivity of an antenna.
This is usually what is meant when the directivity is indicated to be a single number;
in any event, the intended meaning of the term is usually clear from context.  

%%%%%%%%%%%%%%%%%%%%%%%%%%%%%%%%%%%%%%%%%%%%%%%
%ORB EXAMPLE
~\\ \begin{example} 
\begin{mdframed}[backgroundcolor=brown!20,linecolor=brown!20]\raggedright
Directivity of the electrically-short dipole.
\index{electrically-short dipole (ESD)}
\\~\\
An electrically-short dipole (ESD) consists of a straight wire having length $L\ll\lambda/2$.
What is the directivity of the ESD?

{\bf Solution.}
The field radiated by an ESD is derived in Section~\ref{m0198_Radiation_from_an_Electrically-Short_Dipole}.
In that section, we find that the electric field intensity in the far field of a $\hat{\bf z}$-oriented ESD located at the origin is:
%
\begin{equation}
\widetilde{\bf E}({\bf r}) \approx  \hat{\bf \theta} j \eta \frac{I_0\cdot\beta L}{8\pi} ~\left(\sin\theta\right) ~\frac{e^{-j\beta r}}{r}
\label{m0203_eE}
\end{equation}
%
where $I_0$ represents the magnitude and phase of the current applied to the terminals, $\eta$ is the wave impedance of the medium, and $\beta=2\pi/\lambda$.
In Section~\ref{m0207_Total_Power_Radiated_by_an_Electrically-Short_Dipole}, we find that the power density of this field is:
%
\begin{equation}
S({\bf r}) \approx \eta \frac{\left|I_0\right|^2\left(\beta L\right)^2}{128\pi^2} ~\left(\sin\theta\right)^2 ~\frac{1}{r^2}
\label{m0203_eE2}
\end{equation}
%
and we subsequently find that the total power radiated is:
%
\begin{equation}
P_{rad} \approx \eta \frac{\left|I_0\right|^2\left(\beta L\right)^2}{48\pi}
\label{m0203_ePT}
\end{equation}
%
The average power density $S_{ave}$ is simply the total power divided by the area of a sphere centered on the ESD. Let us place this sphere at distance $r$, with $r\gg L$ and $r\gg \lambda$ as required for the validity of Equations~\ref{m0203_eE} and \ref{m0203_eE2}.  Then:
%
\begin{equation}
S_{ave} = \frac{P_{rad}}{4\pi r^2} \approx \eta \frac{\left|I_0\right|^2\left(\beta L\right)^2}{192\pi^2 r^2}
\end{equation}
%
Finally the directivity is determined by applying the definition:
%
\begin{flalign}
D(\theta,\phi) &\triangleq \frac{S({\bf r})}{S_{ave}(r)} \\
               &\approx \underline{1.5 \left(\sin\theta\right)^2}
\end{flalign}

The maximum directivity occurs in the $\theta=\pi/2$ plane.
Therefore, the maximum directivity is \underline{1.5}, meaning the maximum power density is 1.5 times greater than the power density averaged over all directions.

\end{mdframed}
% note figures will need to go between "\end{mdframed}" and "\end{example}"
\end{example}
%%%%%%%%%%%%%%%%%%%%%%%%%%%%%%%%%%%%%%%%%%%%%%%

Since directivity is a unitless ratio, it is common to express it in decibels.  
For example, the maximum directivity of the ESD in the preceding example is $10\log_{10} 1.5 \cong 1.76$~dB.
(Note ``$10\log_{10}$'' here since directivity is the ratio of power-like quantities.)

{\bf Gain.}
The gain $G(\theta,\phi)$ of an antenna is its directivity modified to account for loss within the antenna.
Specifically:
%
\begin{equation}
G(\theta,\phi) \triangleq \frac{S({\bf r})~\mbox{for actual antenna}}{S_{ave}(r)~\mbox{for identical but lossless antenna}}
\end{equation}
%
In this equation, the numerator is the actual power density radiated by the antenna, which is less than the nominal power density due to losses within the antenna. 
The denominator is the average power density for an antenna which is identical, but lossless.
Since the actual antenna radiates less power than an identical but lossless version of the same antenna, gain in any particular direction is always less than directivity in that direction.
Therefore, an equivalent definition of antenna gain is
%
\begin{equation}
G(\theta,\phi) \triangleq e_{rad}D(\theta,\phi)
\end{equation}
%
where $e_{rad}$ is the radiation efficiency 
\index{radiation efficiency}
of the antenna (Section~\ref{m0202_Equivalent_Circuit_Model_for_Transmission}).

%%%%%%%%%%%%%%%%%%%%%%%%%%%%%%%%%%%%%%%%%%%%%%%
%ORB HIGHLIGHT
\begin{mdframed}[backgroundcolor=yellow!20,nobreak=true] 
Gain is directivity times radiation efficiency; that is, directivity modified to account for loss within the antenna.  
\end{mdframed}
%%%%%%%%%%%%%%%%%%%%%%%%%%%%%%%%%%%%%%%%%%%%%%%

{\bf The receive case.}
To conclude this section, we make one additional point about directivity, which applies equally to gain.
The preceding discussion has presumed an antenna which is radiating; i.e., transmitting.  
Directivity can also be defined for the receive case, in which it quantifies the effectiveness of the antenna in converting power in an incident wave to power in a load attached to the antenna.
Receive directivity is formally introduced in Section~\ref{m0218_Effective_Aperture} (``Effective Aperture'').
When receive directivity is defined as specified in Section~\ref{m0218_Effective_Aperture}, it is equal to transmit directivity as defined in this section. 
Thus, it is commonly said that the directivity of an antenna is the same for receive and transmit.  

% Suggestion: Perhaps add that the power received by an antenna from a plane waves from one direction and summed over two polarizations divided by the integrated power received for many plane waves arriving from a full sphere and two polarizations is equal to the directivity, which is defined in transmit mode.

%%%%%%%%%%%%%%%%%%%%%
 
{\bf Additional Reading:}
\begin{itemize}
\item ``\href{https://en.wikipedia.org/wiki/Directivity}{Directivity}'' on Wikipedia.
\end{itemize}

%%%%%%%%%%%%%%%%%%%%%%

\index{directivity|)}

